\clearpage
\begin{examplebox}
\textbf{例1：横向开口，直接验证倒斜率关系。}

抛物线 $y^2=4x$，取 $A=(1,2)$、$B=(4,-4)$，$O=(0,0)$。先核对两点都在曲线上：$2^2=4\cdot1$，$(-4)^2=4\cdot4$。
\[
k_{OA}=\frac21=2,\qquad
k_{OB}=\frac{-4}{4}=-1,\qquad
k_{AB}=\frac{-4-2}{4-1}=-2.
\]
于是
\[
\boxed{\frac12+\frac1{-1}=-\frac12=\frac1{-2}}.
\]
\textbf{反向应用：}不需要计算 $A,B$ 的具体坐标。若已知 $k_{OA}=2$、$k_{OB}=6$，则
\[
\frac1{k_{AB}}=\frac12+\frac16=\frac23
\quad\Longrightarrow\quad\boxed{k_{AB}=\frac32}.
\]
\end{examplebox}

\clearpage
\begin{examplebox}
\textbf{例2：纵向开口，开口向下也成立。}
\begin{center}% Exact example 2: p=-2, A=(2,-1), B=(-4,-4), k_AB=1/2.
\begin{tikzpicture}[font=\small,>=stealth]
\draw[gray!45,rounded corners=3pt] (0,0) rectangle (14.2,9.6);
\node at (7.1,9.05) {开口向下，仍然是三条边的普通斜率关系};
\begin{scope}[shift={(5.0,7.0)},scale=.9]
\draw[->,gray] (-4.7,0)--(3.2,0) node[above] {$x$};
\draw[->,gray] (0,-5.0)--(0,.8) node[left] {$y$};
\draw[teal!65!black,thick,samples=110,domain=-4.5:2.8,variable=\t] plot (\t,{-\t*\t/4});
\fill[blue!4] (0,0)--(2,-1)--(-4,-4)--cycle;
\draw[blue!75!black,very thick] (0,0)--(2,-1);
\draw[orange!85!black,very thick] (0,0)--(-4,-4);
\draw[green!45!black,very thick] (2,-1)--(-4,-4);
\fill (0,0) circle(2pt) node[above left=8pt] {$O$};
\fill (2,-1) circle(2pt) node[above right=8pt] {$A$};
\fill (-4,-4) circle(2pt) node[below left=7pt] {$B$};
\end{scope}
\node[anchor=west,align=left] at (8.9,5.7)
 {$x^2=-4y,\quad p=-2$\\[12pt]
 $A(2,-1)$\\[12pt]$B(-4,-4)$\\[15pt]
 $\textcolor{blue!75!black}{k_{OA}=-\frac12}$\\[12pt]
 $\textcolor{orange!85!black}{k_{OB}=1}$\\[12pt]
 $\textcolor{green!45!black}{k_{AB}=\frac12}$};
\node at (7.1,.7) {$\displaystyle -\frac12+1=\frac12$；$p$ 的正负改变开口方向，公式仍然成立。};
\end{tikzpicture}
\end{center}

设 $x^2=-4y$，即 $p=-2$。取 $A=(2,-1)$、$B=(-4,-4)$，它们均在抛物线上。于是
\[
k_{OA}=\frac{-1}{2}=-\frac12,\qquad
k_{OB}=\frac{-4}{-4}=1.
\]
按照 C097，无须先写弦方程，直接得到
\[
\boxed{k_{AB}=-\frac12+1=\frac12}.
\]
核验：
\[
k_{AB}=\frac{-4-(-1)}{-4-2}
=\frac{-3}{-6}=\frac12.
\]
\textbf{要点：}公式对所有 $p\ne0$ 成立，不要求抛物线一定向右或向上开口。
\end{examplebox}

\clearpage
\begin{examplebox}
\textbf{例3：对称弦——一个不能乱写，另一个完全能用。}
\begin{center}% p=1, free coordinates +1,-1, reflected with identical scale.
\begin{tikzpicture}[font=\small,>=stealth]
\draw[gray!45,rounded corners=3pt] (0,0) rectangle (14.2,10.0);
\node at (7.1,9.45) {对称的两个端点：竖直弦与水平弦的斜率含义不同};
\node at (3.55,8.45) {$y^2=2x$};
\node at (10.65,8.45) {$x^2=2y$};
\begin{scope}[shift={(2.6,5.1)},scale=1.25]
\draw[->,gray] (-1.5,0)--(2.1,0) node[above] {$x$};
\draw[->,gray] (0,-1.5)--(0,2.1) node[left] {$y$};
\draw[teal!65!black,thick,samples=90,domain=-1.4:1.8,variable=\t] plot ({\t*\t/2},\t);
\draw[blue!75!black,thick] (0,0)--(.5,1);
\draw[orange!85!black,thick] (0,0)--(.5,-1);
\draw[green!45!black,very thick] (.5,1)--(.5,-1);
\fill (0,0) circle(2pt) node[below left=7pt] {$O$};
\fill (.5,1) circle(2pt) node[right=9pt] {$A$};
\fill (.5,-1) circle(2pt) node[right=9pt] {$B$};
\end{scope}
\begin{scope}[shift={(9.7,5.1)},scale=1.25]
\draw[->,gray] (-1.5,0)--(2.1,0) node[above] {$x$};
\draw[->,gray] (0,-1.5)--(0,2.1) node[left] {$y$};
\draw[teal!65!black,thick,samples=90,domain=-1.4:1.8,variable=\t] plot (\t,{\t*\t/2});
\draw[blue!75!black,thick] (0,0)--(1,.5);
\draw[orange!85!black,thick] (0,0)--(-1,.5);
\draw[green!45!black,very thick] (-1,.5)--(1,.5);
\fill (0,0) circle(2pt) node[below left=7pt] {$O$};
\fill (1,.5) circle(2pt) node[above=14pt] {$A$};
\fill (-1,.5) circle(2pt) node[above=14pt] {$B$};
\end{scope}
\node at (3.55,2.6) {$AB:\ x=\frac12$，$k_{AB}$ 不存在};
\node at (10.65,2.6) {$AB:\ y=\frac12$，$k_{AB}=0$};
\node at (3.55,1.55) {$\displaystyle\frac{x_B-x_A}{y_B-y_A}=0$};
\node at (10.65,1.55) {$\displaystyle k_{OA}+k_{OB}=\frac12-\frac12=0$};
\node at (7.1,.55) {同一尺度下，横向模型的竖直弦关于 $y=x$ 翻转后成为纵向模型的水平弦。};
\end{tikzpicture}
\end{center}

\textbf{横向：}对 $y^2=2x$，取
\[
A=\left(\frac12,1\right),\qquad B=\left(\frac12,-1\right).
\]
此时 $k_{OA}=2$，$k_{OB}=-2$，故
\[
\frac1{k_{OA}}+\frac1{k_{OB}}=\frac12-\frac12=0.
\]
但 $AB:\ x=1/2$ 为\textbf{竖直直线}，$k_{AB}$ 不存在，因此不允许写 $1/k_{AB}=0$。应使用加强式
\[
\frac{x_B-x_A}{y_B-y_A}=\frac0{-2}=0.
\]

\medskip
\textbf{纵向：}对 $x^2=2y$，取 $A=(1,1/2), B=(-1,1/2)$，则
\[
k_{OA}=\frac12,\qquad k_{OB}=-\frac12,
\qquad\boxed{k_{AB}=0}.
\]
这里 $AB$ 是水平直线，斜率 $0$ \textbf{确实存在}，所以原公式直接成立。
\end{examplebox}

\clearpage
\begin{examplebox}
\textbf{例4：已知三角形两边斜率，求第三边，并判断例外。}

\textbf{(1)} 对 $y^2=6x$，若 $k_{OA}=3$、$k_{OB}=-6$，则
\[
\frac1{k_{AB}}=\frac13-\frac16=\frac16
\quad\Longrightarrow\quad\boxed{k_{AB}=6}.
\]
若改为 $k_{OA}=3$、$k_{OB}=-3$，则两条射线的倒斜率和为 $0$，说明 $AB$ 是竖直弦；此时\textbf{不能}说 $k_{AB}$ 等于无穷大，因为通常的斜率定义中竖直线斜率不存在。

\medskip
\textbf{(2)} 对 $x^2=6y$，若 $k_{OA}=3$、$k_{OB}=-3$，则
\[
\boxed{k_{AB}=3-3=0},
\]
弦 $AB$ 水平，这是完全正常的情况。

\textbf{解题选择：}知道两条射线斜率，优先直接使用公式；知道两点坐标，优先使用平方差；知道弦方程，优先考虑联立和韦达。
\end{examplebox}

\clearpage
\begin{examplebox}
\textbf{例5：顶点不在原点，照样用顶点连线斜率和。}
\begin{center}% Exact translation: Y=X^2 -> y=(x-2)^2+1; equal x/y units.
\begin{tikzpicture}[font=\small,>=stealth]
\draw[gray!45,rounded corners=3pt] (0,0) rectangle (14.2,10.0);
\node at (7.1,9.45) {只平移图形，三条边的方向都不变};
\node at (3.55,8.55) {$Y=X^2$};
\node at (10.65,8.55) {$y=(x-2)^2+1$};
\begin{scope}[shift={(3.2,3.8)},scale=1.15]
\draw[->,gray] (-2.1,0)--(2.1,0) node[above] {$X$};
\draw[->,gray] (0,-.75)--(0,3.4) node[left] {$Y$};
\draw[teal!65!black,thick,samples=80,domain=-1.7:1.7,variable=\t] plot (\t,{\t*\t});
\draw[blue!75!black,very thick] (0,0)--(1,1);
\draw[orange!85!black,very thick] (0,0)--(-1.5,2.25);
\draw[green!45!black,very thick] (1,1)--(-1.5,2.25);
\fill (0,0) circle(2pt) node[below left=8pt] {$V$};
\fill (1,1) circle(2pt) node[right=12pt] {$A$};
\fill (-1.5,2.25) circle(2pt) node[right=14pt] {$B$};
\end{scope}
\begin{scope}[shift={(8.1,2.65)},scale=1.15]
\draw[->,gray] (-.25,0)--(4.5,0) node[above] {$x$};
\draw[->,gray] (0,-.3)--(0,4.4) node[left] {$y$};
\draw[teal!65!black,thick,samples=80,domain=.3:3.7,variable=\t] plot (\t,{(\t-2)*(\t-2)+1});
\draw[blue!75!black,very thick] (2,1)--(3,2);
\draw[orange!85!black,very thick] (2,1)--(.5,3.25);
\draw[green!45!black,very thick] (3,2)--(.5,3.25);
\draw[gray!55,dashed] (2,0)--(2,1)--(0,1);
\fill (2,1) circle(2pt) node[below right=10pt] {$V(2,1)$};
\fill (3,2) circle(2pt) node[right=12pt] {$A$};
\fill (.5,3.25) circle(2pt) node[right=14pt] {$B$};
\end{scope}
\node at (3.55,1.8) {$k_{VA}=1,\quad k_{VB}=-\frac32$};
\node at (10.65,1.8) {$k_{VA}=1,\quad k_{VB}=-\frac32$};
\node at (7.1,.65) {两图同尺度：横、纵坐标差保持不变，故 $k_{AB}=1-\frac32=-\frac12$。};
\end{tikzpicture}
\end{center}
在 $y=(x-2)^2+1$ 上，取顶点 $V(2,1)$、$A(3,2)$、$B(1/2,13/4)$。代入可核对两端点在曲线上，且均不同于顶点。
\[
k_{VA}=\frac{2-1}{3-2}=1,\qquad
k_{VB}=\frac{13/4-1}{1/2-2}=-\frac32.
\]
直接应用推广，得 $\boxed{k_{AB}=1-3/2=-1/2}$。
两点式回验：
\[
\frac{13/4-2}{1/2-3}=\frac{5/4}{-5/2}=-\frac12.
\]
\textbf{识别依据：}必须连接实际顶点 $V$，不能仍把坐标原点当成顶点。
\end{examplebox}

\clearpage
\begin{examplebox}
\textbf{例6：只知道两端切线斜率，也能直接求弦斜率。}
\begin{center}% x^2=2y, p=1; A=(-1,1/2), B=(2,2); slopes -1,2,1/2.
\begin{tikzpicture}[font=\small,>=stealth]
\draw[gray!45,rounded corners=3pt] (0,0) rectangle (14.2,10.0);
\node at (7.1,9.45) {纵向：两端切线的斜率取平均，得到弦的斜率};
\begin{scope}[shift={(3.8,3.1)},scale=1.15]
\draw[->,gray] (-2.4,0)--(3.4,0) node[above] {$x$};
\draw[->,gray] (0,-1.3)--(0,4.3) node[left] {$y$};
\draw[teal!65!black,thick,samples=100,domain=-2.2:2.8,variable=\t] plot (\t,{\t*\t/2});
\draw[blue!75!black,dashed,thick] (0,0)--(-1,.5);
\draw[orange!85!black,dashed,thick] (0,0)--(2,2);
\draw[blue!75!black,very thick] (-2.2,1.7)--(.6,-1.1);
\draw[orange!85!black,very thick] (.6,-.8)--(2.8,3.6);
\draw[green!45!black,very thick] (-1,.5)--(2,2);
\fill (0,0) circle(2pt) node[below left=10pt] {$O$};
\fill (-1,.5) circle(2pt) node[above=14pt] {$A$};
\fill (2,2) circle(2pt) node[right=14pt] {$B$};
\node[blue!75!black] at (-1.6,2.1) {$t_A$};
\node[orange!85!black] at (3.25,3.25) {$t_B$};
\end{scope}
\node[anchor=west,align=left] at (8.9,5.7)
 {$x^2=2y,\quad p=1$\\[13pt]
 $A(-1,\frac12),\ B(2,2)$\\[15pt]
 $k_A=-1,\quad k_B=2$\\[15pt]
 $k_{AB}=\dfrac{-1+2}{2}=\dfrac12$};
\node at (7.1,1.3) {虚线是顶点连线，实线 $t_A,t_B$ 是切线；二者不能混淆。};
\node at (7.1,.6) {$k_{OA}=-\frac12,\ k_{OB}=1$，切线斜率分别是顶点连线斜率的 $2$ 倍。};
\end{tikzpicture}
\end{center}
\textbf{纵向：}对 $x^2=2y$，两端切线斜率为 $-1,2$，则
\[
\boxed{k_{AB}=\frac{-1+2}{2}=\frac12}.
\]
回验：由切线斜率 $k=x$，两点为 $A(-1,1/2),B(2,2)$；两点式得 $(2-1/2)/(2+1)=1/2$。

\textbf{横向变式：}对 $y^2=2x$，若切线斜率为 $-1,1/2$，则
\[
\frac1{k_{AB}}=\frac12\left(-1+2\right)=\frac12,
\qquad\boxed{k_{AB}=2}.
\]
回验两点为 $(1/2,-1),(2,2)$，弦斜率 $3/(3/2)=2$。
\textbf{关键：}先辨开口，再辨题目给的是顶点连线还是切线。
\end{examplebox}

\clearpage
\begin{examplebox}
\textbf{例7：同一条抛物线，两种焦点弦的长度倒数和。}
\begin{center}% p=2; y^2=4x. Left chord x=1; right A=(4,4),B=(1/4,-1).
\begin{tikzpicture}[font=\small,>=stealth]
\draw[gray!45,rounded corners=3pt] (0,0) rectangle (14.2,11.0);
\node at (7.1,10.45) {同一条抛物线，焦点弦方向改变，长度倒数和不变};
\node at (3.55,9.55) {竖直焦点弦};
\node at (10.65,9.55) {倾斜焦点弦};
\begin{scope}[shift={(1.6,4.9)},scale=.85]
\draw[->,gray] (-.5,0)--(5.2,0) node[above] {$x$};
\draw[->,gray] (0,-2.9)--(0,4.5) node[left] {$y$};
\draw[teal!65!black,thick,samples=100,domain=-2.7:4.25,variable=\t] plot ({\t*\t/4},\t);
\draw[blue!75!black,very thick] (1,0)--(1,2);
\draw[orange!85!black,very thick] (1,0)--(1,-2);
\fill (1,2) circle(2pt) node[right=12pt] {$A$};
\fill (1,-2) circle(2pt) node[right=12pt] {$B$};
\fill (1,0) circle(2pt) node[below right=12pt] {$F$};
\end{scope}
\begin{scope}[shift={(8.7,4.9)},scale=.85]
\draw[->,gray] (-.5,0)--(5.2,0) node[above] {$x$};
\draw[->,gray] (0,-2.9)--(0,4.5) node[left] {$y$};
\draw[teal!65!black,thick,samples=100,domain=-2.7:4.25,variable=\t] plot ({\t*\t/4},\t);
\draw[blue!75!black,very thick] (1,0)--(4,4);
\draw[orange!85!black,very thick] (1,0)--(.25,-1);
\fill (4,4) circle(2pt) node[above left=9pt] {$A$};
\fill (.25,-1) circle(2pt) node[right=14pt] {$B$};
\fill (1,0) circle(2pt) node[below right=12pt] {$F$};
\end{scope}
\node at (3.55,2.1) {$FA=FB=2$};
\node at (10.65,2.1) {$FA=5,\quad FB=\frac54$};
\node at (3.55,1.2) {$\frac12+\frac12=1$};
\node at (10.65,1.2) {$\frac15+\frac45=1$};
\node at (7.1,.6) {两图同尺度：$y^2=4x,\ p=2,\ F(1,0)$；相加的是长度倒数。};
\end{tikzpicture}
\end{center}
对 $y^2=4x$，$p=2$、$F(1,0)$，所以任何有限两点焦点弦都满足
\[
\boxed{\frac1{FA}+\frac1{FB}=1}.
\]
\textbf{最简单的验证：}$x=1$ 的端点为 $(1,\pm2)$，两段均长 $2$。

\textbf{一般位置的数值回验：}$A(4,4),B(1/4,-1)$ 的坐标满足曲线方程；连线为 $y=\frac43(x-1)$，确过焦点。距离分别为 $5,5/4$，故倒数和为 $1/5+4/5=1$。
这里 $AB$ 竖直也能使用\textbf{长度}倒数和，不能把这个结论误换成弦的\textbf{斜率}倒数。
\end{examplebox}

\clearpage
\begin{examplebox}
\textbf{例8：乘积除以和，先看分母是否为零。}
\begin{center}% p=1, free coordinates +1,-1, reflected with identical scale.
\begin{tikzpicture}[font=\small,>=stealth]
\draw[gray!45,rounded corners=3pt] (0,0) rectangle (14.2,10.0);
\node at (7.1,9.45) {对称的两个端点：竖直弦与水平弦的斜率含义不同};
\node at (3.55,8.45) {$y^2=2x$};
\node at (10.65,8.45) {$x^2=2y$};
\begin{scope}[shift={(2.6,5.1)},scale=1.25]
\draw[->,gray] (-1.5,0)--(2.1,0) node[above] {$x$};
\draw[->,gray] (0,-1.5)--(0,2.1) node[left] {$y$};
\draw[teal!65!black,thick,samples=90,domain=-1.4:1.8,variable=\t] plot ({\t*\t/2},\t);
\draw[blue!75!black,thick] (0,0)--(.5,1);
\draw[orange!85!black,thick] (0,0)--(.5,-1);
\draw[green!45!black,very thick] (.5,1)--(.5,-1);
\fill (0,0) circle(2pt) node[below left=7pt] {$O$};
\fill (.5,1) circle(2pt) node[right=9pt] {$A$};
\fill (.5,-1) circle(2pt) node[right=9pt] {$B$};
\end{scope}
\begin{scope}[shift={(9.7,5.1)},scale=1.25]
\draw[->,gray] (-1.5,0)--(2.1,0) node[above] {$x$};
\draw[->,gray] (0,-1.5)--(0,2.1) node[left] {$y$};
\draw[teal!65!black,thick,samples=90,domain=-1.4:1.8,variable=\t] plot (\t,{\t*\t/2});
\draw[blue!75!black,thick] (0,0)--(1,.5);
\draw[orange!85!black,thick] (0,0)--(-1,.5);
\draw[green!45!black,very thick] (-1,.5)--(1,.5);
\fill (0,0) circle(2pt) node[below left=7pt] {$O$};
\fill (1,.5) circle(2pt) node[above=14pt] {$A$};
\fill (-1,.5) circle(2pt) node[above=14pt] {$B$};
\end{scope}
\node at (3.55,2.6) {$AB:\ x=\frac12$，$k_{AB}$ 不存在};
\node at (10.65,2.6) {$AB:\ y=\frac12$，$k_{AB}=0$};
\node at (3.55,1.55) {$\displaystyle\frac{x_B-x_A}{y_B-y_A}=0$};
\node at (10.65,1.55) {$\displaystyle k_{OA}+k_{OB}=\frac12-\frac12=0$};
\node at (7.1,.55) {同一尺度下，横向模型的竖直弦关于 $y=x$ 翻转后成为纵向模型的水平弦。};
\end{tikzpicture}
\end{center}
对 $y^2=2x$，给出两组顶点连线斜率。

\textbf{正例：}$k_{OA}=2,k_{OB}=6$。和为 $8\ne0$，故
\[
\boxed{k_{AB}=\frac{2\cdot6}{2+6}=\frac32}.
\]
由 $y=2x,y=6x$ 与曲线的另一交点得 $A(1/2,1),B(1/18,1/3)$，回验弦斜率为 $(-2/3)/(-4/9)=3/2$。

\textbf{边界例：}$k_{OA}=2,k_{OB}=-2$。分母为 $0$，不可计算 $-4/0$；两点为 $(1/2,\pm1)$，弦是 $x=1/2$，普通斜率不存在。
若换成纵向 $x^2=2y$，顶点连线斜率为 $2,-2$ 时，弦斜率直接为 $0$，不需要除以斜率和。
\end{examplebox}

\clearpage
\begin{examplebox}
\textbf{短例9：另外两种“倒数相加”，各用各的条件。}

\textbf{(1) 椭圆焦点弦。}椭圆 $x^2/25+y^2/9=1$ 的半轴为 $a=5,b=3$，右焦点为 $F(4,0)$，所以长度倒数和为
\[
\boxed{\frac{2a}{b^2}=\frac{10}{9}}.
\]
沿长轴的两端点为 $(5,0),(-5,0)$，距离为 $1,9$，倒数和 $1+1/9=10/9$。
垂直长轴的焦点弦两端点为 $(4,\pm9/5)$，代入椭圆方程得 $16/25+9/25=1$；倒数和也是 $5/9+5/9=10/9$。

\medskip
\textbf{(2) $3,4,5$ 直角三角形。}面积为 $3\cdot4/2=6$，斜边为 $5$，斜边上的高为 $h=2\cdot6/5=12/5$。于是
\[
\frac1{h^2}=\frac{25}{144},\qquad
\frac1{3^2}+\frac1{4^2}=\frac{16+9}{144}=\frac{25}{144}.
\]
长度平方的倒数对应勾股定理，不能改写成 $1/h=1/3+1/4$；实际 $5/12\ne7/12$。
\end{examplebox}
